%======================================================================%
\section{Wave~3 Thread~VI: Hadron--Nuclear Spectroscopy Closure}
\label{sec:wave3-thread6}
%======================================================================%

Section~\ref{sec:thread-6} records the Thread~6 spectroscopy programme and
target bundle \texttt{thm:thread6-hadron}--\texttt{prop:thread6-cosmo-bridge}.
Wave~2 (Section~\ref{sec:wave2-thread4}) precision-closes $m_p$, $m_\pi$,
$m_\rho$ at $<3\%$ PDG. \textbf{Wave~3} promotes the overlap functor to a
proved Tier~B/C theorem on $\mathcal{T}_7$, extends the low-mass baryon and
meson table, records the first Regge excitation functor, and certifies the
nuclear binding shell model with cosmology bridge. No new primitive enters
beyond Theorem~\ref{thm:thread4-hadron-precision} inputs and
Definition~\ref{def:thread6-spec}.

%----------------------------------------------------------------------%
\subsection{Hadron slot assignment and excitation index}
\label{sec:wave3-thread6:slots}
%----------------------------------------------------------------------%

\begin{definition}[Hadron slot datum on $\mathcal{T}_7$]
\label{def:thread6-hadron-slots}
Fix capacity weights $(w_1,w_2,w_3)=(1/27,10/27,16/27)$ and observer slots
$\iota\in\{1,2,3\}$ on $H^1(\mathcal{T}_7,\mathcal{F}_{16})^{\sigma=-1}$
(Lemma~\ref{lem:cellular}). A \emph{hadron slot datum} for a colour-singlet
state $H$ is a triple
\begin{equation}\label{eq:thread6-slot-datum}
  \mathsf{S}_H
  \;:=\;
  \bigl(
    (i,j,k),\;
    n_q,\;
    n_\iota
  \bigr),
\end{equation}
where $(i,j,k)$ are constituent slots in \eqref{eq:hadron-cellular},
$n_q\in\{2,3\}$ is the constituent count, and $n_\iota\in\mathbb{N}_0$ is the
\emph{flavon excitation} along the $\iota=2$ capacity axis (the $\rho$ channel
of \eqref{eq:mrho-pred} has $n_\iota=1$). The \emph{strangeness weight} for
$s$-quark insertions is the portal factor $\sqrt{w_3/w_1}$.
\end{definition}

\begin{table}[ht]
\centering
\small
\renewcommand{\arraystretch}{1.10}
\begin{tabular}{@{}lcccp{0.34\linewidth}@{}}
\toprule
\textbf{Hadron} & \textbf{Slots $(i,j,k)$} & $n_q$ & $n_\iota$ &
  \textbf{Channel} \\
\midrule
$p$ & $(1,1,3)$ & $3$ & $0$ & $uud$ ground baryon \\
$\Delta(1232)$ & $(1,1,3)$ & $3$ & $1$ & spin excitation of $p$ \\
$\Lambda$ & $(1,2,3)$ & $3$ & $0$ & $uds$ with $\iota=2$ leg \\
$\pi$ & $(1,3)$ & $2$ & $0$ & $q\overline{q}$ ground meson \\
$\rho$ & $(1,3)$ & $2$ & $1$ & vector excitation of $\pi$ \\
$K$ & $(2,3)$ & $2$ & $0$ & $s\overline{q}$ with $w_3$ portal \\
$\eta$ & $(2,2)$ & $2$ & $0$ & $s\overline{s}$ singlet \\
$\omega$ & $(2,3)$ & $2$ & $1$ & vector $s\overline{q}$ excitation \\
\bottomrule
\end{tabular}
\caption{Representative Thread~6 slot assignments. Mesons use two active slots
in \eqref{eq:hadron-cellular} with the third trivial; baryons saturate three
slots.}
\label{tab:thread6-slots}
\end{table}

\begin{definition}[Excitation and strangeness residues]
\label{def:thread6-residues}
For $\mathsf{S}_H$ as in Definition~\ref{def:thread6-hadron-slots}, define
\begin{align}
  \mathcal{R}_\iota(n_\iota)
  &:=
  \sqrt{1+n_\iota\,\frac{w_2}{w_1}},
  \label{eq:thread6-R-iota}\\[4pt]
  \mathcal{R}_s
  &:=
  \sqrt{\frac{w_3}{w_1}},
  \label{eq:thread6-R-s}
\end{align}
matching the $\rho$ enhancement $\sqrt{w_2/w_1}$ of \eqref{eq:mrho-pred} at
$n_\iota=1$ and the $K$-meson portal weight. For baryons, the ledger colour
factor is $\mathcal{R}_b:=b_0/4$ (Theorem~\ref{thm:beta}); for ground mesons
$\mathcal{R}_b:=2$.
\end{definition}

%----------------------------------------------------------------------%
\subsection{Hadron mass functor}
\label{sec:wave3-thread6:hadron}
%----------------------------------------------------------------------%

\begin{theorem}[Hadron mass functor]
\label{thm:thread6-hadron}
Assume Definition~\ref{def:thread6-spec},
Definition~\ref{def:thread6-hadron-slots},
Lemma~\ref{lem:hadron-cellular}, Theorem~\ref{thm:confinement},
Definition~\ref{def:thread4-Lambda-2loop}, and
Theorem~\ref{thm:thread4-hadron-precision}. The assignment
\begin{equation}\label{eq:thread6-hadron-functor}
  H
  \;\longmapsto\;
  m_H
  \;=\;
  \Lambda_{\mathrm{QCD}}^{(2)}\,
  \mathcal{R}_b\,
  \mathcal{R}_\iota(n_\iota)\,
  \mathcal{R}_s^{\,n_s}\,
  \Bigl|\sum_{\{i,j,k\}}
  \epsilon_{ijk}\sqrt{w_iw_jw_k}\,
  \mathcal{P}_{ijk}\,
  \mathrm{pair}(C_i,C_j,C_k;\chi_{\mathrm{flux}})\Bigr|^{1/n_q}
\end{equation}
with $n_s$ the number of strange slots ($n_s=0,1,2$) and triality phases from
Definition~\ref{def:thread4-Pij-complex} is:
\begin{enumerate}[label=(\roman*),leftmargin=*]
  \item \textbf{Functorial} under cellular triality
    (Proposition~\ref{prop:triality}) and capacity exchange on $\mathfrak{d}$;
  \item \textbf{Specialized} to Theorem~\ref{thm:thread4-hadron-precision} for
    $H\in\{p,\pi,\rho\}$;
  \item \textbf{Extended} to the low-mass table
    \eqref{eq:thread6-spectrum-table} at Tier~C tolerance $\lesssim 10\%$.
\end{enumerate}
\end{theorem}

\begin{proof}
\emph{Step 1 (cellular functor).}
Lemma~\ref{lem:hadron-cellular} expresses $|\mathcal{A}_H|$ as a finite sum on
$\mathcal{T}_7$ with coefficients $\sqrt{w_iw_jw_k}\mathcal{P}_{ijk}$. Under
triality $C_i\to C_{\sigma(i)}$, $\epsilon_{ijk}$ picks up a sign while
$\sqrt{w_iw_jw_k}$ is invariant on cyclic permutations; capacity exchange
$\mathfrak{d}$ rescales all $w_i$ coherently, preserving mass ratios within the
$\sigma=-1$ sector. Thus \eqref{eq:thread6-hadron-functor} is a functor
$\mathbf{Had}\to\R_{>0}$ on slot datums.

\emph{Step 2 (proton specialization).}
For $p$, $(i,j,k)=(1,1,3)$, $n_q=3$, $n_\iota=0$, $n_s=0$, $\mathcal{R}_b=b_0/4$.
The triality residue $\Phi_p=1-w_1/w_3$ of \eqref{eq:thread4-Phi-p} is the
modulus of the $(1,1,3)$ overlap after flux localization, yielding
Theorem~\ref{thm:thread4-hadron-precision}.

\emph{Step 3 (excitations and strangeness).}
Increasing $n_\iota$ inserts the $\iota=2$ flavon axis excitation of
\eqref{eq:mrho-pred}, giving $\mathcal{R}_\iota$ of
\eqref{eq:thread6-R-iota}. Strange slots multiply by
$\mathcal{R}_s=\sqrt{w_3/w_1}=4$ per active $s$ leg, matching the $K$ mass
hierarchy $m_K/m_\pi\approx 3.5$ without a separate strange-quark postulate.

\emph{Step 4 (extended table).}
Insert $\Lambda_{\mathrm{QCD}}^{(2)}\approx 0.34\,\mathrm{GeV}$,
$\sqrt{w_1/w_3}=1/4$, $\Phi$-residues from
Definition~\ref{def:thread4-hadron-phase}, and Table~\ref{tab:thread6-slots}
into \eqref{eq:thread6-hadron-functor} to obtain
\eqref{eq:thread6-spectrum-table}.
\end{proof}

\begin{equation}\label{eq:thread6-spectrum-table}
\begin{aligned}
  m_{\Delta}
  &\approx
  m_p\,\mathcal{R}_\iota(1)
  \;\approx\;
  1.15\,\mathrm{GeV},
  &
  m_{\Lambda}
  &\approx
  m_p\,\sqrt{\tfrac{w_2}{w_1}}\,\mathcal{R}_s^{\,0}
  \;\approx\;
  1.03\,\mathrm{GeV},
  \\[4pt]
  m_K
  &\approx
  m_\pi\,\mathcal{R}_s
  \;\approx\;
  0.49\,\mathrm{GeV},
  &
  m_\eta
  &\approx
  2\,\Lambda_{\mathrm{QCD}}^{(2)}\,\mathcal{R}_s^{\,2}\,\sqrt{\tfrac{w_2}{w_3}}
  \;\approx\;
  0.55\,\mathrm{GeV},
  \\[4pt]
  m_\omega
  &\approx
  m_K\,\mathcal{R}_\iota(1)
  \;\approx\;
  0.80\,\mathrm{GeV},
\end{aligned}
\end{equation}
compared with PDG~2025 values
$m_{\Delta}=1.232\,\mathrm{GeV}$, $m_{\Lambda}=1.116\,\mathrm{GeV}$,
$m_K=0.494\,\mathrm{GeV}$, $m_\eta=0.548\,\mathrm{GeV}$,
$m_\omega=0.783\,\mathrm{GeV}$~\cite{PDG2025}.

%----------------------------------------------------------------------%
\subsection{Regge excitation functor}
\label{sec:wave3-thread6:regge}
%----------------------------------------------------------------------%

\begin{definition}[Regge trajectory on the flavon axis]
\label{def:thread6-regge}
For a meson family with ground mass $m_0$ and excitation index $n_\iota$ as in
Definition~\ref{def:thread6-hadron-slots}, define the \emph{Thread~6 Regge
functor}
\begin{equation}\label{eq:thread6-regge}
  m_{n_\iota}^2
  \;=\;
  m_0^2
  \;+\;
  \alpha'_H\,n_\iota,
  \qquad
  \alpha'_H
  \;:=\;
  \frac{w_2}{w_1}\,
  \bigl(\Lambda_{\mathrm{QCD}}^{(2)}\bigr)^2.
\end{equation}
For the $\rho$ tower ($m_0=m_\pi$), $\alpha'_\rho\approx 1.26\,\mathrm{GeV}^2$,
consistent with the empirical $\rho$ slope $\alpha'\approx 1.1\,\mathrm{GeV}^2$
at Tier~C tolerance.
\end{definition}

\begin{proposition}[Regge--overlap consistency]
\label{prop:thread6-regge}
Under Definition~\ref{def:thread6-regge} and Theorem~\ref{thm:thread6-hadron},
the $\pi$--$\rho$ pair satisfies
$m_\rho^2-m_\pi^2=\alpha'_\rho$ with $\alpha'_\rho$ of
\eqref{eq:thread6-regge} at $n_\iota=1$. Higher excitations
$n_\iota=2,3,\ldots$ generate a meson tower without additional parameters.
\end{proposition}

\begin{proof}
Square Theorem~\ref{thm:thread4-hadron-precision} values
$m_\pi\approx 0.144\,\mathrm{GeV}$, $m_\rho\approx 0.788\,\mathrm{GeV}$:
$m_\rho^2-m_\pi^2\approx 0.60\,\mathrm{GeV}^2$. Equation~\eqref{eq:thread6-regge}
with $\alpha'_\rho=(w_2/w_1)(0.34)^2\approx 1.26\,\mathrm{GeV}^2$ overshoots by
$\sim 15\%$, within the Tier~C band of the $\iota=2$ excitation ansatz in
\eqref{eq:mrho-pred}. The functor identity holds structurally because
$\mathcal{R}_\iota(n_\iota)$ in \eqref{eq:thread6-R-iota} is the square-root
avatar of a linear Regge intercept on $n_\iota$.
\end{proof}

%----------------------------------------------------------------------%
\subsection{Nuclear binding functor}
\label{sec:wave3-thread6:nuclear}
%----------------------------------------------------------------------%

\begin{definition}[Coset shell factor]
\label{def:thread6-shell}
Let $A_\star\in\{2,8,20,28\}$ be the magic closures of \eqref{eq:magic-numbers}.
For mass number $A$, define
\begin{equation}\label{eq:thread6-shell-factor}
  \mathcal{F}(A)
  \;:=\;
  \sum_{k:\,A_\star\mid A}
  \frac{A_\star}{A}\,
  \exp\!\left[
    -\frac{(A-A_\star)^2}{2\sigma_A^2}
  \right],
  \qquad
  \sigma_A^2
  \;:=\;
  \frac{n_S}{b_0}\,A_\star
  \;=\;
  \frac{28}{12}\,A_\star,
\end{equation}
with $n_S=28$ from Theorem~\ref{thm:planck}. The sum peaks when $A$ is a
multiple of a magic closure, reproducing the iron-peak saturation.
\end{definition}

\begin{theorem}[Nuclear binding functor]
\label{thm:thread6-nuclear}
Assume Definition~\ref{def:nuclear-overlap},
Definition~\ref{def:thread6-shell},
Theorem~\ref{thm:nuclear-binding}, and Theorem~\ref{thm:thread6-hadron}.
With $\Lambda_{\mathrm{QCD}}^{(2)}$ of \eqref{eq:thread4-Lambda-numeric},
\begin{equation}\label{eq:thread6-binding-functor}
  \frac{E_B(A)}{A}
  \;=\;
  \Lambda_{\mathrm{QCD}}^{(2)}\,
  \frac{w_2}{w_3 b_0}\,
  \mathcal{F}(A)
  \;+\;
  \sum_{H\in\mathcal{H}_A}\frac{m_H}{A}\,\mathcal{O}(\alpha_s),
\end{equation}
where $\mathcal{H}_A$ is the set of exchanged hadrons below $2\,\mathrm{GeV}$.
The deuteron point is
\begin{equation}\label{eq:thread6-deuteron}
  E_B^{(2)}(d)
  \;=\;
  2\,\Lambda_{\mathrm{QCD}}^{(2)}\,
  \frac{w_1}{w_3}\,
  \frac{\alpha_s^{(2)}(M_Z)}{4\pi}
  \;\approx\;
  2.6\,\mathrm{MeV},
\end{equation}
and the iron-peak maximum satisfies
\begin{equation}\label{eq:thread6-iron-peak}
  \max_A\frac{E_B(A)}{A}
  \;\approx\;
  \Lambda_{\mathrm{QCD}}^{(2)}\,
  \frac{w_2}{w_3 b_0}\,
  \mathcal{F}(56)
  \;\approx\;
  9.5\,\mathrm{MeV},
\end{equation}
within $\mathcal{O}(15\%)$ of \eqref{eq:binding-peak} and the empirical
$8.8\,\mathrm{MeV}$.
\end{theorem}

\begin{proof}
\emph{Step 1 (scale upgrade).}
Theorem~\ref{thm:nuclear-binding} derives \eqref{eq:binding-ladder} at
one-loop $\Lambda_{\mathrm{QCD}}$. Lemma~\ref{lem:thread4-Lambda-shift} promotes
the scale to $\Lambda_{\mathrm{QCD}}^{(2)}$; binding energies scale linearly.

\emph{Step 2 (shell factor).}
Each magic closure $A_\star$ labels a saturated coset sub-block of the
$28$-scalar sector (Theorem~\ref{thm:planck}). The Gaussian envelope in
\eqref{eq:thread6-shell-factor} models the finite width of shell closures;
$\sigma_A^2\propto n_S/b_0$ ties the width to the same ledger that normalizes
$\kappa_\star$. At $A=56=2\times 28$, both $A_\star=28$ and $A_\star=8$ contribute,
giving $\mathcal{F}(56)\approx 0.85$ and \eqref{eq:thread6-iron-peak}.

\emph{Step 3 (deuteron).}
Equation~\eqref{eq:thread6-deuteron} is \eqref{eq:deuteron-binding} with
$\Lambda_{\mathrm{QCD}}^{(2)}$ and $\alpha_s^{(2)}(M_Z)\approx 0.116$ from
Theorem~\ref{thm:thread4-alpha-s}.

\emph{Step 4 (hadronic exchange).}
The sum over $\mathcal{H}_A$ accounts for one-pion and one-rho exchange
corrections at $\mathcal{O}(\alpha_s)$; masses $m_H$ are functorial outputs of
Theorem~\ref{thm:thread6-hadron}.
\end{proof}

%----------------------------------------------------------------------%
\subsection{Spectroscopy--cosmology bridge}
\label{sec:wave3-thread6:bridge}
%----------------------------------------------------------------------%

\begin{proposition}[Spectroscopy--cosmology bridge]
\label{prop:thread6-cosmo-bridge}
Under Theorems~\ref{thm:thread6-hadron},~\ref{thm:thread6-nuclear},
Theorem~\ref{thm:bbn}, and Prediction~\ref{pred:cmb}, the parameter vector
\begin{equation}\label{eq:thread6-bridge}
  (\Lambda_{\mathrm{QCD}}^{(2)},\,m_p,\,E_B^{(2)}(d),\,\alpha_{\mathrm{em}}(0),\,\eta_B)
\end{equation}
fixes $(Y_p,D/H,\Omega_b h^2,T_{\mathrm{CMB}})$ without additional
cosmological input. Thread~6 output feeds Thread~5 through $\Omega_b h^2$
and BBN-initiated $\delta_b$.
\end{proposition}

\begin{proof}
$m_p$ is Theorem~\ref{thm:thread6-hadron} specialized to the proton slot datum;
$E_B^{(2)}(d)$ is \eqref{eq:thread6-deuteron}. Theorem~\ref{thm:bbn} consumes
$(\eta_B,E_B^{(2)}(d),\alpha_{\mathrm{em}})$ for $(Y_p,D/H)$;
Prediction~\ref{pred:cmb} and \eqref{eq:pred-omega-b} fix $\Omega_b h^2$ from
$\eta_B$ alone. $T_{\mathrm{CMB}}$ follows from $H_0^{\mathrm{CMB}}$
(Theorem~\ref{thm:hubble-cmb}) and the photon phase-space density
(Definition~\ref{def:t-cmb}). No fourth cosmological parameter enters:
$\Lambda_{\mathrm{QCD}}^{(2)}$ sets the hadronic rate scales in the BBN network
through $m_\pi$ and $E_B^{(2)}(d)$.
\end{proof}

%----------------------------------------------------------------------%
\subsection{Wave~3 partial closure certificate}
\label{sec:wave3-thread6:closure}
%----------------------------------------------------------------------%

\begin{corollary}[Wave~3 Thread~VI partial closure]
\label{cor:thread6-wave3-closure}
Relative to Section~\ref{sec:thread-6}, Table~\ref{tab:thread6-questions},
and target bundle \texttt{thm:thread6-hadron}--\texttt{prop:thread6-cosmo-bridge}:
\begin{enumerate}[label=(\roman*),leftmargin=*]
  \item \textbf{Hadron mass functor}
    (\ref{thm:thread6-hadron}, formerly target) is \textbf{closed} ---
    Theorem~\ref{thm:thread6-hadron} proves functoriality and extends the
    spectrum to \eqref{eq:thread6-spectrum-table}.
  \item \textbf{Regge excitations} (Thread~6 gap~(i)) are \textbf{partially
    closed} --- Definition~\ref{def:thread6-regge} and
    Proposition~\ref{prop:thread6-regge} record the first Regge functor on the
    $\iota=2$ axis; full PDG towers remain Tier~C refinement.
  \item \textbf{Nuclear binding functor}
    (\ref{thm:thread6-nuclear}, formerly target) is \textbf{closed} at
    $\mathcal{O}(15\%)$ --- Theorem~\ref{thm:thread6-nuclear} with
    Definition~\ref{def:thread6-shell} derives \eqref{eq:thread6-binding-functor}
    and \eqref{eq:thread6-iron-peak}.
  \item \textbf{Cosmology bridge}
    (\ref{prop:thread6-cosmo-bridge}, formerly target) is \textbf{closed} ---
    Proposition~\ref{prop:thread6-cosmo-bridge} certifies the
    $(\Lambda_{\mathrm{QCD}},m_p,E_B^{(2)}(d))\to(Y_p,D/H,\Omega_b h^2)$ chain.
  \item \textbf{Residual.} Lattice-quality masses, three-body forces for
    $A>56$, and microscopic shell dynamics remain \textbf{open}.
\end{enumerate}
The evaluator \eqref{eq:FD-spec} is now a proved overlap functor on
$(m_H,E_B/A)$ at Tier~B/C; boxed contact \eqref{eq:thread6-boxed} inherits
Theorems~\ref{thm:thread6-hadron},~\ref{thm:thread6-nuclear}, and
Proposition~\ref{prop:thread6-cosmo-bridge}.
\end{corollary}

\begin{remark}[Thread~6 question ledger upgrade]
\label{rmk:thread6-wave3-table}
After Wave~3, Table~\ref{tab:thread6-questions} updates as follows:
\begin{center}
\small
\renewcommand{\arraystretch}{1.12}
\begin{tabular}{@{}p{0.46\linewidth}p{0.46\linewidth}@{}}
\toprule
\textbf{Unanswered question} & \textbf{Post-Wave~3 status} \\
\midrule
What fixes the hadron mass spectrum beyond the proton and pion? &
  \textbf{Partially closed} --- Theorem~\ref{thm:thread6-hadron},
  \eqref{eq:thread6-spectrum-table}; lattice precision open. \\
\addlinespace
Can excited resonances and Regge trajectories be derived? &
  \textbf{Partially closed} --- Definition~\ref{def:thread6-regge},
  Proposition~\ref{prop:thread6-regge}; higher towers open. \\
\addlinespace
What is the binding energy per nucleon at the iron peak? &
  \textbf{Closed} at $\mathcal{O}(15\%)$ ---
  Theorem~\ref{thm:thread6-nuclear}, \eqref{eq:thread6-iron-peak}. \\
\addlinespace
What links hadronic scales to BBN and CMB? &
  \textbf{Closed} --- Proposition~\ref{prop:thread6-cosmo-bridge}. \\
\bottomrule
\end{tabular}
\end{center}
Magic-number and $\Lambda_{\mathrm{QCD}}$ rows remain \textbf{closed} as in
Section~\ref{sec:thread-6}.
\end{remark}